\begin{theorem}[Group left absorption forces the right factor to the unit]
\label{thm:group-left-absorb-right-factor-unit}
For a history-level group operation with associativity, left and right
identities, multiplication congruence, and left inverses, if $ab\hsame a$,
then $b\hsame e$.
\end{theorem}
\begin{proof}
Chain $ab$ through the assumption to $a$, then backward through the right
identity $ae\hsame a$. Left cancellation of the common factor $a$ yields
$b\hsame e$.
\end{proof}
\leanchecked{BEDC.Derived.GroupUp.group\_left\_absorb\_right\_factor\_unit}

\begin{theorem}[Group left multiplication cancellation from empty unit]
\label{thm:group-left-mul-cancel-empty-unit-iff}
For a BEDC history-level group operation with empty unit $\emp$,
associativity, left empty-unit law, multiplication congruence, and left
inverses, left multiplication by a common factor preserves and reflects
$\hsame$:
$$
\hsame(ab,ac) \Longleftrightarrow \hsame(b,c).
$$
\end{theorem}
\begin{proof}
The forward direction is left cancellation for the common factor $a$. The
reverse direction transports $b\hsame c$ through multiplication congruence with
the fixed left factor $a$.
\end{proof}
\leanchecked{BEDC.Derived.GroupUp.group\_left\_mul\_cancel\_from\_empty\_unit\_iff}

\begin{theorem}[Group right cancellation]
\label{thm:group-right-cancellation}
For a history-level group operation with associativity, right identity,
multiplication congruence, and right inverses, equality of \(ba\) and \(ca\)
under \(\hsame\) implies equality of \(b\) and \(c\) under \(\hsame\).
\end{theorem}
\begin{proof}
Chain \(b\) through the right identity, replace that identity by
\(aa^{-1}\), reassociate to \((ba)a^{-1}\), use multiplication congruence on
the assumed product equality, reassociate back to \(c(aa^{-1})\), collapse the
right inverse, and finish by the right identity on \(c\).
\end{proof}
\leanchecked{BEDC.Derived.GroupUp.group\_right\_cancel}

\begin{theorem}[Group right multiplication cancellation from empty unit]
\label{thm:group-right-mul-cancel-empty-unit-iff}
For a BEDC history-level group operation with empty unit $\emp$,
associativity, right empty-unit law, multiplication congruence, and right
inverses, right multiplication by a common factor preserves and reflects
$\hsame$:
$$
\hsame(ba,ca) \Longleftrightarrow \hsame(b,c).
$$
\end{theorem}
\begin{proof}
The forward direction is right cancellation for the common factor $a$. The
reverse direction transports $b\hsame c$ through multiplication congruence with
the fixed right factor $a$.
\end{proof}
\leanchecked{BEDC.Derived.GroupUp.group\_right\_mul\_cancel\_from\_empty\_unit\_iff}

\begin{theorem}[Group right absorption forces the left factor to the unit]
\label{thm:group-right-absorb-left-factor-unit}
For a history-level group operation with associativity, left and right
identities, multiplication congruence, and right inverses, if $ba\hsame a$,
then $b\hsame e$.
\end{theorem}
\begin{proof}
Chain $ba$ through the assumption to $a$, then backward through the left
identity $ea\hsame a$. Right cancellation of the common factor $a$ yields
$b\hsame e$.
\end{proof}
\leanchecked{BEDC.Derived.GroupUp.group\_right\_absorb\_left\_factor\_unit}

\begin{theorem}[Group two-sided cancellation]
\label{thm:group-two-sided-cancellation}
For a history-level group operation with associativity, left and right
identities, multiplication congruence, and two-sided inverses, equality of
$(ab)c$ and $(ad)c$ under $\hsame$ implies equality of $b$ and $d$ under
$\hsame$.
\end{theorem}
\begin{proof}
First apply right cancellation to the common final factor $c$, reducing the
displayed equality to $\hsame(ab,ad)$. Then apply left cancellation to the
common initial factor $a$, yielding $\hsame(b,d)$.
\end{proof}
\leanchecked{BEDC.Derived.GroupUp.group\_two\_sided\_cancel}

\begin{definition}[Group ledger policy]
\label{def:group-ledger-policy}
Record the monoid ledger plus every inverse operation applied.
\end{definition}

\begin{definition}[\(\GroupUp\) certificate obligations]
\label{def:group-certificate-obligations}
The carrier, source, pattern, classifier, stability, and ledger data listed above are the obligations required for the interface name \(\GroupUp\).
\end{definition}

\paragraph{Primitive boundary.}
\(\GroupUp\) is not a BEDC primitive. It is admitted only when those obligations are supplied.

\begin{proposition}[Group forgets to Monoid certificate]
\label{prop:group-forgets-monoid-certificate}
Let $C$ carry a $\GroupUp(C)$ certificate with multiplication $\cdot$, identity $e$, inverse operation ${}^{-1}$, and classifier $\sim_C$. If the inverse operation, inverse closure, inverse congruence, left and right inverse rows, and ledger entries whose expression provenance uses inverse formation are omitted, then the retained carrier, source, monoid pattern, classifier, monoid stability fields, and monoid ledger form a $\MonoidUp(C)$ certificate. The carrier $C$, multiplication $\cdot$, identity $e$, and classifier $\sim_C$ are unchanged.
\end{proposition}
\begin{proof}
By \autoref{def:group-certificate-obligations}, the given certificate supplies carrier, source, pattern, classifier, stability, and ledger data for $\GroupUp(C)$. The construction keeps exactly the coordinates whose operations and endpoint claims use only the monoid part of the group structure.

By \autoref{def:group-carrier}, a group carrier is a monoid setup extended with an inverse operation. Omitting the inverse operation leaves the same carrier $C$, the same classifier $\sim_C$, the same multiplication $\cdot$, and the same identity $e$. These are precisely the carrier data listed in \autoref{def:monoid-carrier}.

By \autoref{def:group-source-specification}, the source is a certified monoid carrier plus inverse data. Removing the inverse component leaves the certified monoid carrier required by \autoref{def:monoid-source-specification}. No source endpoint, carrier witness, multiplication operation, identity element, or classifier is replaced.

By \autoref{def:group-pattern-specification}, the group pattern is monoid expression formation extended by inverse formation. Restricting to expressions with no inverse node leaves expression formation from $e$, carrier witnesses, and binary multiplication, which is exactly the monoid pattern of \autoref{def:monoid-pattern-specification}.

By \autoref{def:group-classifier-specification}, the group classifier is the monoid carrier classifier together with inverse congruence. The retained relation is therefore the same carrier classifier required by \autoref{def:monoid-classifier-specification}; the omitted inverse congruence row is not a $\MonoidUp$ obligation.

By \autoref{def:group-stability-certificate}, the group stability certificate consists of all monoid stability fields together with inverse congruence, left inverse law, right inverse law, and closure under inverse. The first component is exactly the list in \autoref{def:monoid-stability-certificate}: classifier reflexivity, classifier transitivity, associativity under classifier, left identity, right identity, and multiplication congruence. The inverse-specific rows mention ${}^{-1}$ and are not required for $\MonoidUp$.

Finally, \autoref{def:group-ledger-policy} records the monoid ledger plus every inverse operation applied. Omitting entries justified by inverse formation leaves the carrier witnesses and multiplication-expression provenance described by \autoref{def:monoid-ledger-policy}. Hence the retained carrier, source, pattern, classifier, stability, and ledger data satisfy the obligations named in \autoref{def:monoid-certificate-obligations}, with $C$, $\cdot$, $e$, and $\sim_C$ copied from the original $\GroupUp(C)$ certificate.
\end{proof}

\begin{definition}[Opposite group certificate data]
\label{def:opposite-group-certificate-data}
Let a $\GroupUp(C)$ certificate be presented by a carrier predicate $\mathcal C$, a classifier $\sim_C$, multiplication $\cdot_C$, identity $e_C$, and inverse operation $i_C$. Define the opposite multiplication by
$$
a\star_C b:=b\cdot_C a.
$$
The opposite group data $C^{\mathrm{op}}$ have the same carrier predicate $\mathcal C$, the same classifier $\sim_C$, the same identity $e_C$, the same inverse operation $i_C$, and multiplication $\star_C$.
\end{definition}

\begin{theorem}[Opposite group certificate]
\label{thm:opposite-group-certificate}
If $(\mathcal C,\sim_C,\cdot_C,e_C,i_C)$ satisfies the $\GroupUp$ certificate obligations of \autoref{def:group-certificate-obligations}, then the data of \autoref{def:opposite-group-certificate-data} also satisfy the $\GroupUp$ certificate obligations. Thus the same carrier, classifier, identity, and inverse operation carry a group certificate with multiplication $\star_C$.
\end{theorem}
\leanchecked{BEDC.Derived.GroupUp.GroupSingletonHistory\_opposite\_laws}
\begin{proof}
Unpack the group obligations through \autoref{def:group-certificate-obligations}. The monoid reduct supplied by \autoref{prop:group-forgets-monoid-certificate} has carrier $\mathcal C$, classifier $\sim_C$, multiplication $\cdot_C$, and identity $e_C$. Applying \autoref{thm:opposite-monoid-certificate} to that reduct gives the carrier, source, monoid pattern, classifier, monoid stability rows, and monoid ledger for the multiplication $\star_C$ with the same identity $e_C$.

These rows supply exactly the monoid part demanded by \autoref{def:group-source-specification}, \autoref{def:group-pattern-specification}, and the first component of \autoref{def:group-stability-certificate}. The carrier data of \autoref{def:group-carrier} are also present: the carrier, classifier, and identity are unchanged, the multiplication is $\star_C$, and the inverse operation is still $i_C$.

It remains to check the inverse-specific group rows. Inverse closure is copied from the original group certificate because the inverse operation and carrier predicate are unchanged: from $\mathcal C(a)$ the original row gives $\mathcal C(i_C(a))$. Inverse congruence is copied for the same reason: from $a\sim_C b$ the original row gives $i_C(a)\sim_C i_C(b)$, and the classifier is still $\sim_C$.

The two inverse laws exchange their handedness under the reversed multiplication. The opposite left-inverse row at a carried element $a$ unfolds to
$$
i_C(a)\star_C a=a\cdot_C i_C(a)\sim_C e_C,
$$
which is the original right-inverse row. The opposite right-inverse row unfolds to
$$
a\star_C i_C(a)=i_C(a)\cdot_C a\sim_C e_C,
$$
which is the original left-inverse row.

The classifier row of \autoref{def:group-classifier-specification} is the same carrier classifier equipped with the inverse congruence just copied. For the pattern and ledger, a $\star_C$ group-expression tree is interpreted as a $\cdot_C$ group-expression tree by exchanging the two children at each binary multiplication node while leaving carrier leaves, the identity leaf, and inverse nodes unchanged. The opposite monoid certificate supplies the binary-node source and ledger reading; the inverse nodes use the original inverse source and ledger entries of \autoref{def:group-ledger-policy}. Thus the source, pattern, classifier, stability, and ledger rows form the $\GroupUp$ certificate data required by \autoref{def:typed-naming-certificate} and \autoref{def:group-certificate-obligations}.
\end{proof}
